\centerline{\Huge \bf  Задачи к экзамену по алгебре.}

\section{\Large \bf Дроби}

\problem{1} Что больше: бутылка на троих или две на семерых?

\textit{Решение:} $1/3 = 7/21, 2/7 = 6/21,$ т. е. $1/3 > 2/7.$

\problem{2} Как отрезать от шнура в 2/3 м кусок в 1/2 м, не имея метра?

\textit{Решение:} Кусок в 1/2 м составляет три четверти от всего шнура:
\[ \dfrac{3}{4} \cdot \dfrac{2}{3} = \dfrac{1}{2}\]
так что от шнура надо отрезать четверть (два раза сложив его пополам).


\problem{3} Как разделить семь апельсинов между 12 лицеистами, если ни один апельсин нельзя резать более чем на пять частей?

\textit{Решение:} Надо сначала разделить между учениками три апельсина(каждому достанется по четверти апельсина), а затем разделить между ними оставшиеся четыре апельсина (каждому $–$ по трети).

\problem{4}
Упростите выражение:
\[
\frac{(x-a)(x-b)}{(c-a)(c-b)} + \frac{(x-a)(x-c)}{(b-a)(b-c)} + \frac{(x-b)(x-c)}{(a-b)(a-c)}
\]



\textit{Решение:} Сложим две первые дроби, приведя их к общему знаменателю \((c-a)(c-b)(b-a)\) (числитель и знаменатель первой дроби умножаем на \(b-a\), числитель и знаменатель второй -- на \(c-a\); используем, что $b - c = -(c - b):$


\[
\frac{(x-a)(x-b)}{(c-a)(c-b)} + \frac{(x-a)(x-c)}{(b-a)(b-c)} =
\frac{(x-a)(x-b)(b-a)-(x-a)(x-c)(c-a)}{(c-a)(c-b)(b-a)} =
\]

\[
\frac{(x-a)[(x-b)(b-a)-(x-c)(c-a)]}{(c-a)(c-b)(b-a)} =
\frac{(x-a)[xb - xa - b^2 + ab - xc + xa + c^2 - ac]}{(c-a)(c-b)(b-a)} =
\]

\[
\frac{(x-a)[xb - xc + a(b - c) - (b - c)(b + c)]}{(c-a)(c-b)(b-a)} =
\frac{(x-a)(b-c)(x + a - b - c)}{(c-a)(c-b)(b-a)}.
\]

Сократив на \((c-b) = -(b-c)\), получаем

\[
\frac{(x-a)(b+c-a-x)}{(c-a)(b-a)}.
\]

\newpage

\hspace{3mm}

Прибавим теперь третью дробь (знаменатели у них одинаковы):

\[
\frac{(x-a)(b+c-a-x)}{(c-a)(b-a)} + \frac{(x-b)(x-c)}{(a-b)(a-c)} =
\]

\[
\frac{xb + xc - xa - x^2 - ab + ac + a^2 + ax + x^2 - xb - xc + bc}{(c-a)(b-a)} =
\]

\[
\frac{a^2 + bc - ab - ac}{(c-a)(b-a)} = \frac{a(a-b-c) + c(a-b) + b(a-c)}{(c-a)(b-a)} = 1.
\]

Итак, мы доказали тождество

\[
\frac{(x-a)(x-b)}{(c-a)(c-b)} + \frac{(x-a)(x-c)}{(b-a)(b-c)} + \frac{(x-b)(x-c)}{(a-b)(a-c)} = 1.
\]



\section{\Large \bf Делимость}


\problem{1}
Докажите, что для всех натуральных n число, записываемое $3^n$ единицами, делится на $3^n$.

\problem{2}
Докажите признаки делимости на: 

a) 4 и 5

б) 3 и 9

в) 11

г)$^*$ $2^n$ и $5^n$

\problem{3}
Рассмотрим число N, записанное в десятичной системе счисления. Найдём сумму цифр этого числа, потом сложим цифры, которыми записана сумма и т.д. Будем продолжать этот процесс, пока в конце концов не получим однозначное число, которое называют \textit{цифровым корнем} числа N.

а) Докажите, что цифровой корень сравним с N по модулю 9.

б) Выведите явную формулу

\problem{4}
Натуральные числа $a, b, c, d$ таковы, что  $ab = cd.$  Докажите, что найдутся такие натуральные $u, v, w, z,$ что $a = uv,  b = wz,  c = uw,  d = vz.$

\textit{Решение:} Пусть $u = НОД(a, c).$  Тогда  $a = uv,  c = uw,$  где $v$ и $w$ взаимно просты. Из данного равенства следует, что  $vb = wd.$  Значит, $d$ делится на $v:  d = vz$. Подставляя в равенство и сокращая, получаем  $b = wz.$

\newpage

\hspace{3mm}

\section{Исследование квадратного уравнения.}

\problem{1}
Докажите, что квадрат имеет максимальную площадь среди всех прямоугольников данного периметра.

\textit{Решение:} Пусть нам даны прямоугольники с периметром равным $P$. Тогда если одна из сторон равна $x$, то другая равна $\dfrac{P}{2} - x$. Найдём площадь этого прямоугольника. 

\[x\cdot (\dfrac{P}{2} - x) = -x^2 +\dfrac{P}{2} \cdot x.\]

Ветви параболы направлены вниз, значит в вершине квадратный трёхчлен принимает максимальное значение. Значит при $x = \dfrac{P}{4}$ площадь максимальна. А это и значит, что наш прямоугольник квадрат.

\problem{2}
Уравнение \(x^2 + px + q = 0\) имеет корни \(x_1\) и \(x_2\). Выразите \((x_1 - x_2)^2\) через \(p\) и \(q\).

\textit{Первое решение.} \((x_1 - x_2)^2 = x_1^2 - 2x_1x_2 + x_2^2 = x_1^2 + x_2^2 - 4x_1x_2 = (x_1 + x_2)^2 - 4x_1x_2 = p^2 - 4q.\)

\textit{Второе решение.} \(x_1 - x_2\) $-$ разность между корнями; глядя на формулу, видим, что она равна \(2\sqrt{D}\), так что \((x_1 - x_2)^2 = 4D = 4 \left( \frac{p^2}{4} - q \right) = p^2 - 4q.\)

\problem{3}
\textit{(Теорема Виета для кубического уравнения)} Уравнение
\[
x^3 + px^2 + qx + r = 0
\]
имеет три (различных) корня \(\alpha, \beta, \gamma\). Доказать, что
\[
\alpha + \beta + \gamma = -p,
\]
\[
\alpha \beta + \alpha \gamma + \beta \gamma = q,
\]
\[
\alpha \beta \gamma = -r.
\]


\problem{4$^{*}$} \textit{(Возвратное уравнение.)} Найдите корни уравнения $x^4 + 3x^3 + 2x^2 + 3x + 1 = 0$

\textit{Решение:} Поделим уравнение $x^2$, благо $x=0$ не является корнем. Получим 

\[ x^2 + 3x + 2 + 3\dfrac{1}{x} + \dfrac{1}{x^2} = (x^2 + 2 + \dfrac{1}{x^2}) + 3 (x+\dfrac{1}{x}) = (x+\dfrac{1}{x})^2 + 3(x+\dfrac{1}{x}) = 0\].

Обозначим $x + \dfrac{1}{x}$ за $y$ и решим новое квадратное уравнение.

\[y^2 + 3y = 0 \to y = 0; \ y = 3 \to  
\begin{cases}
   x+\dfrac{1}{x} = 0 &\text{нет решений.}\\
   \\
   x + \dfrac{1}{x} = 3 & x = \dfrac{-3 \pm \sqrt{5}}{2}
 \end{cases}
 \]

 \section{Функции и их графики}

 \problem{1} Найдите координаты вершины (нижней или верхней точки) графика $y = a(x+m)^2 + n$

 \textit{Решение:} Дабы минимизировать правое выражение, стоит обнулить скобку $(x+m)^2$ подставив туда $-m$. Стало быть вершина находится в точке $(-m;n)$

 \problem{2}
 На рисунке изображен график функции  $y = (a^2 - 1)(x^2 - 1) + (a - 1)(x-1).$ Найдите координаты точки А.

 \textit{Решение:} Так как графиком является прямая, то заданная функция $-$ линейная. Следовательно,  $a^2 - 1 = 0  \longleftrightarrow  a = \pm 1.$  Рассмотрим два случая.
 
  1)  $a = 1$.  Тогда  $y = 0,$  что не соответствует заданному графику.
  
  2)  $a = -1.$  Тогда  $y = - 2x + 2,$  что соответствует заданному \\графику. Осталось подставить  $x = 0.$ и получить ответ $A(0, 2)$

\begin{wrapfigure}{r}{0.2\linewidth}
    \vspace{-30mm}
    \includegraphics[width=\linewidth]{ЭлМШ 2023-2024/VIII блок/Киви/lesson_8/image_1.png}
\end{wrapfigure}

\problem{3} Алина забыла продержурить после урока и на доске остался график функции  $y = \dfrac{k}{x}$  и пять прямых, параллельных прямой  $y = kx  (k \neq 0)$. Найдите произведение абсцисс всех точек пересечения прямых с гиперболой.

\textit{Решение:} Прямая, параллельная прямой  $y = kx$,  имеет уравнение  $y = kx + b$.  Абсциссами точек её пересечения с гиперболой являются оба корня уравнения   $\dfrac{k}{x} = kx + b$  $\longleftrightarrow$   $kx^2 + bx - k = 0.$   Произведение корней этого уравнения равно –1. Перемножив пять таких произведений, получаем ответ.

\textit{Замечания:}

1. Каждое из указанных квадратных уравнений имеет два действительных корня. Геометрически это как раз означает, что любая прямая, параллельная прямой $y = kx$,  пересекает гиперболу в двух точках.

2. Аналогично можно доказать более общий факт – произведение абсцисс точек пересечения прямой и гиперболы  $y = \dfrac{k}{x}$  зависит только от $k$ и угла наклона прямой.

\newpage

\hspace{3mm}

\problem{4} На улице $5$ домов. Каждый день почтальон идёт на почту, берёт там письма для жителей одного дома и разносит их. Затем он возвращается на почту, берёт письма для жителей другого дома и снова их разносит. И так он обходит все дома. В каком месте нужно построить почту, чтобы почтальону пришлось проходить наименьшее расстояние? Улицу можно считать отрезком прямой.

\textit{Решение:} Введём координатную прямую, на которой расположим дома с координатами $x_1, x_2, x_3, x_4, x_5$ упорядоченными по возрастанию. Пусть $t \ -$ координата почтового отделения на числовой прямой.

\begin{figure}[h]
\centering
\includegraphics[width=0.8\linewidth]{ЭлМШ 2023-2024/VIII блок/Киви/lesson_8/image_2.png}
\end{figure}

 Почтальон проходит от почты до каждого дома дважды: туда и обратно, то есть общее расстояние, пройденное почтальоном равно
$2S = 2|x_1 - t| + 2|x_2 - t| + 2|x_3 - t| + 2|x_4 - t| + 2|x_5 - t|.$   Нужно сделать 2S, а значит и S, как можно меньше. Рассмотрим функцию $S(t).$ Графиком является ломаная, точки излома в каждой из точек $x_1, x_2, x_3, x_4, x_5$ (см. рисунок).

\begin{figure}[h]
\centering
\includegraphics[width=0.5\linewidth]{ЭлМШ 2023-2024/VIII блок/Киви/lesson_8/image_3.png}
\end{figure}

У самого левого звена угловой коэффициент $-5$, а если двигаться вправо, то при переходе через каждую точку $x_i$ угловой коэффициент увеличивается на $2$. Таким образом, до точки $x_3$ функция $S(t)$ убывает, а после этой точки возрастает (на рисунке масштаб по оси ординат уменьшен для удобства).
  Значит, наименьшее значение достигается в точке $t = x_3$.  Здание почты нужно строить в этой точке.  

